% English translation of questions/5780/semA-moedA/partB-q2.tex (metadata is taken from the Hebrew file).

\begin{question}
Let $f(x)$ be differentiable on the interval $(a,b)$. Which of the following must be true?
\begin{enumerate}
  \item[(a)] $f'(x)$ is continuous on $(a,b)$
  \item[(b)] $f(x)$ is monotone on $(a,b)$
  \item[(c)] $f(x)+\sin(x)$ is continuous on $(a,b)$
  \item[(d)] $f(x)$ is monotone increasing on $(a,b)$
  \item[(e)] there exists $c\in(a,b)$ such that $f'(c)=0$
\end{enumerate}
\end{question}

\begin{hint}
What does differentiability at a point imply?
\end{hint}

\begin{solution}
\textbf{The correct answer: (c).}

A function differentiable at a point is continuous there, hence $f$ is continuous on $(a,b)$. Also $\sin x$ is continuous, and a sum of continuous functions is continuous; hence $f(x)+\sin(x)$ is continuous on $(a,b)$.

The other answers need not hold:
\begin{itemize}
  \item (a) is false: the derivative of a differentiable function need not be continuous. The classical example: $f(x)=x^2\sin\frac1x$ for $x\neq0$ and $f(0)=0$, on the interval $(-1,1)$. It is differentiable at every point ($f'(0)=\lim_{h\to0}h\sin\frac1h=0$), but $f'(x)=2x\sin\frac1x-\cos\frac1x$ for $x\ne0$, and $f'$ has no limit at $0$.
  \item (b), (d) are false: for example $f(x)=x^2$ on the interval $(-1,1)$ is differentiable but not monotone.
  \item (e) is false: for example $f(x)=x$ is differentiable and $f'(x)=1\ne0$ for every $x$. (Rolle's theorem also requires continuity on a closed interval and $f(a)=f(b)$.)
\end{itemize}
\end{solution}
